\begin{proposition}[7.3.12]
\label{IV.2.s7.3.12-fr}
Sea $\boldsymbol P$ una propiedad del primer o del segundo tipo definida
a partir de una propiedad $\boldsymbol R$ que satisfaga las condiciones
$(\mathrm R_{\mathrm I})$ y $(\mathrm R_{\mathrm{II}})$ \textup{(resp.}
$(\mathrm R'_{\mathrm I})$ y $(\mathrm R_{\mathrm{II}})$). Si, para
todo preesquema localmente noetheriano $X$, $U_{\mathbf{R}}(X)$
designa el conjunto de los $x\in X$ para los cuales se cumple
$\boldsymbol R(\mathscr O_x)$, entonces, para todo morfismo regular
\textup{(resp.} todo $\boldsymbol P$-morfismo) $f:X\to Y$ de preesquemas
localmente noetherianos, se tiene
\begin{equation}
U_{\mathbf{R}}(X)=f^{-1}(U_{\mathbf{R}}(Y)).
\tag{7.3.12.1}
\label{IV.2.eq.7.3.12.1-fr}
\end{equation}
\end{proposition}

\begin{proof}
Es una consecuencia inmediata de las definiciones.
\end{proof}
